\documentclass[instructions.tex]{subfiles}
\begin{document}
 
\chapter{Des bons livres.}
 
IL faut lire de bons livres pour s'instruire de la religion, et pour conserver le souvenir de ce qu'on en sait déjà. Il faut en lire encore pour nourrir son âme de pensées et de sentimens salutaires. Mais si l'on veut éviter de graves inconvéniens, il est indispensable d'user de beaucoup de prudence dans le choix des livres; et si l'on désire recueillir un fruit solide de ses lectures, il faut lire avec précaution et avec méthode.
 
\primo Tout homme qui ne connaît pas Dieu, dit saint Jérôme, est un animal. Saint Bernard veut qu'il soit encore de pire condition. Cependant une ignorance qui va jusqu'à ce point n'est pas aussi rare qu'on le pense. saint Paul la reprochait déjà à quelques fidèles de son temps. Il est vrai, l'on n'ignore pas qu'il y a un Dieu : serait-il possible qu'il y eût des athées de conviction? Mais combien ont de Dieu des idées fausses! combien qui rabaissent sa majesté infinie, limitent sa puissance sans bornes, avilissent sa sagesse incompréhensible, nient sa providence qui embrasse l'univers, font injure à sa justice et à sa miséricorde, dont le règne s'étend sur tous les êtres créés! N'est-ce pas ignorer Dieu que de lui attribuer, par exemple, une indifférence oisive pour les choses de la terre? que de prétendre qu'il a livré le monde à un hasard aveugle, ou à un destin inévitable ou irrésistible? que de s'imaginer qu'il se soucie peu qu'on l'aime ou qu'on le haïsse, qu'on le loue ou qu'on le blasphême, qu'on le serve ou qu'on le déshonore, qu'on marche dans les routes du vice, ou qu'on suive les sentiers de la vertu? Est-ce connaître Dieu que de croire qu'il a placé dans certains astres une influence secrète sur le caractère, sur les passions et la destinée des hommes en cette vie; dans des songes fortuits, les pronostics de ce qui doit arriver d'heureux et de funeste; dans certaines rencontres, des annonces de bonheur ou de malheur; dans certains jours, le succès prospère ou contraire d'une entreprise, etc.? A combien de chrétiens ne peut-on pas dire avec l'Apôtre, qu'ils ignorent Dieu, qu'ils le méconnaissent et l'outragent par les idées indignes qu'ils s'en forment\footnote{Ignorantiam Dei quidam habent. 1. Cor. 15. 34.}!
 
On en voit encore qui, exempts de ces excès condamnables, tombent dans d'autres également pernicieux. Leur ignorance en matière de religion va jusqu'à ne savoir point qu'il y a un seul Dieu en trois personnes réellement distinctes; que c'est la seconde personne qui s'est revêtue de notre nature; que cet Homme-Dieu a souffert, a été crucifié, est mort pour nous racheter du péché, de l'esclavage du démon et de la mort éternelle. Ils ne connaissent guère mieux l'immortalité de notre âme, l'éternité des peines pour les méchans, des récompenses pour les bons. A l'égard des sacremens, à peine en savent-ils le nombre et qui les a institués. Si vous leur demandez quelle est la nature de chaque sacrement en particulier, quels en sont les effets, avec quelles dispositions il faut les recevoir, ils ne peuvent vous répondre. Ils ne sont pas plus instruits sur les commandemens. Ils prient sans presque aucune notion sur les dispositions qu'exige la prière, sur les effets qu'elle opère. Ils vont à la messe, sans savoir ce que c'est que la messe; à la table sainte, sans pouvoir dire ce qu'on y reçoit : on en a vu qui croyaient que la sainte hostie contenait le symbole des apôtres.
 
Des chrétiens si peu instruits ont-ils une foi si suffisante? adorent-ils Dieu en esprit et en vérité? sont-ils en état de profiter des moyens que la religion nous offre pour nous aider dans l'affaire périlleuse de notre salut? sont-ils capables de recevoir aucun sacrement? et pourrait-on, sans prévariquer, les admettre à l'office intéressant de parrain et de marraine?
 
Et comment remplissent-ils les principaux devoirs du christianisme, ainsi que les obligations les plus étroites de leur état? S'ils sont chefs de famille, ils élèvent leurs enfans comme des brutes.
 
L'incertitude, le doute, la présomption, l'erreur, le vice; voilà les compagnons presque inséparables de l'ignorance en matière de religion : on pourrait y en ajouter encore d'autres non moins fâcheux.
 
O vous qui pensez si indignement de Dieu et de ses augustes attributs, hâtez-vous d'en chercher des idées plus saines et plus conformes à la vérité. Il s'est révélé lui-même; apprenez dans de bons livres ce qu'il a daigné nous faire connaître de lui. Et vous, qui avilissez le Père des humains par vos conceptions puériles et superstitieuses, étudiez aussi dans des livres choisis la religion, qui n'enseigne que des dogmes dignes du souverain Être, de sa providence soigneuse pour le gouvernement du monde, de sa bonté, de sa justice et de sa miséricorde envers les hommes. Il n'est point de moyen plus efficace pour bannir promptement les ténèbres de l'ignorance, que des lectures lumineuses, appropriées, faites avec de bonnes dispositions. Ceux qui ne savent pas lire doivent assister aux instructions familières de leurs pasteurs, ou se faire instruire en particulier; ils n'ont point d'autres ressources, et rien ne peut les exempter de ce devoir.
 
\secundo Autant la lecture des mauvais livres est pernicieuse par les effets déplorables qu'elle opère sur l'esprit et le cœur de la plupart de ceux qui ont l'imprudence de se la permettre, autant la lecture des bons livres est utile à ceux qui y ont recours avec de sages dispositions.
 
Dieu se servit autrefois de ce moyen pour dessiller les yeux à des idolâtres, pour les tirer des ténèbres profondes dans lesquelles ils étaient plongés dès leur enfance, pour leur faire connaître la beauté de la doctrine chrétienne, leur en inspirer de l'amour, et les engager à embrasser une religion qui parle dignement de la Divinité, qui montre à l'homme sa véritable grandeur, la vraie source de sa dégradation et de sa misère, qui le console dans ses peines, le secourt dans sa faiblesse, l'anime au bien par des promesses magnifiques pour la vie future. Plusieurs reconnurent à ces traits que notre religion sainte avait son origine dans le ciel, et qu'elle était un présent que le Tout-Puissant avait fait aux hommes dans sa miséricorde. Que d'hérétiques aussi ont, par le même moyen, ouvert les yeux à la lumière de la foi, et, renonçant à leurs erreurs, sont rentrés sincèrement dans le sein de l'unité catholique! Que d'hommes, long-temps infatués des vanités du monde, en ont enfin connu le vide et les ont abjurées, afin de servir Dieu loin des dangers, avec plus de liberté et moins d'entraves! Ne sont-ils pas nombreux les exemples de conversion qu'ont donnés de grands pécheurs, lesquels ayant vu, dans d'excellens livres, une vive peinture de l'état affreux de leur âme, des périls dont ils étaient assiégés de toutes parts, des calamités effroyables auxquelles les exposaient, pour l'éternité, les désordres de leur vie passée, se sont convertis, sont revenus franchement à Dieu et ont persévéré avec courage, jusqu'à leur dernière heure, dans les exercices de la pénitence et des autres vertus chrétiennes? N'est-ce pas encore la lecture que nous conseillons, qui a soutenu une multitude d'âmes saintes dans les voies étroites de la justice, les empêchant de se laisser emporter au gré des passions, séduire par les maximes perverses, entraîner par les exemples pernicieux?
 
S'étonnera-t-on de ces effets admirables, quand on considérera qu'un bon livre est, dans la main de la miséricorde céleste qui nous parle en même temps par la grace intérieure, un maître qu'elle nous donne, afin que nous puissions puiser auprès de lui, en tout temps, des leçons pleines de lumières et de sagesse?
 
En effet, ici il nous apprend ce que nous devons croire pour avoir une foi distincte de nos augustes mystères. Là il nous montre le bien qui nous est commandé, et nous porte à l'entreprendre, en même temps qu'il nous suggère les moyens et la manière de le faire; plus loin, c'est un vice qu'il nous dépeint sous de hideuses couleurs, afin de nous en inspirer de l'éloignement; ensuite, descendant avec nous au fond de notre cœur, il en examine tous les plis, et, s'il y trouve quelques vestiges de ce vice, il nous indique ce que nous devons faire pour le détruire en nous; si nous en sommes exempts, il nous prescrit les mesures à prendre pour nous en préserver à jamais. Commençons-nous à nous laisser abattre par une lassitude dangereuse et presque désespérante? il nous console, nous rassure et nous exhorte. Présumons-nous, au contraire, de la miséricorde divine, ou nous reposons-nous sur nos propres forces et sur les victoires déjà remportées sur l'ennemi de notre salut? il nous représente alors avec force notre faiblesse extrême, la multitude des dangers séduisans qui nous environnent, la sévérité rigoureuse des jugemens de Dieu, l'incertitude ou nous nous trouvons de notre état, ne sachant point si nous sommes dignes d'amour ou de haine\footnote{Nescit homo utrùm amore an odio, dignus sit. Eccl. 9. 1.}. Enfin, c'est Dieu qui nous parle par ce bon livre, dit saint Augustin\footnote{Quandò legis, Deus tibi loquitur. Enarr. in Psal. 85. F.}.
 
Concluons donc que nous avons besoin de recourir souvent à ce sage maître. Car nous portons au dedans de nous un foyer de penchans malheureux, qui nous pousse vivement au mal; au dehors, nous vivons au milieu du monde, qui nous met continuellement en péril; le démon nous fait aussi la guerre sans relâche, cherchant les endroits faibles par où il pourra faire entrer la tentation dans notre imagination, et de là dans notre cœur; enfin nos occupations nous dissipent, et en nous dissipant elles interrompent, pour ainsi dire, l'exercice de cette vigilance habituelle, sans laquelle nous avons tout lieu de craindre la surprise. Ainsi peu à peu les impressions salutaires qu'ont produites en nous les instructions que nous avons entendues, s'affaiblissent; les bonnes pensées, les pieux mouvemens, qui ont été le résultat de ces instructions, se présentent plus rarement, et nos résolutions enfin nous échappent. Une lecture réfléchie et fréquente viendrait puissamment à notre secours. On voit des ouvriers et des ouvrières qui tiennent des romans ouverts devant leurs yeux, et y jettent de temps en temps des regards; hélas! dans quel dessein? Pourquoi n'en ferions-nous pas autant, pour nous armer contre les dangers qui nous environnent sans cesse, et pour nous animer à travailler courageusement à notre salut? Du moins imitons ces familles pieuses, où l'on fait une lecture spirituelle, tous les soirs, immédiatement avant de se retirer pour aller prendre son repos.
 
\tertio Mais s'il est utile de lire pour s'instruire et se fortifier, il n'est pas moins nécessaire d'user d'une très-grande prudence dans le choix des livres qu'on se propose de se procurer et de lire. Il y en a beaucoup qui, sous des titres apparens et trompeurs, cachent une doctrine empoisonnée et remplie d'erreurs publiques que l'Église a proscrites.
 
Le moyen de se préserver du danger de la séduction est facile. C'est de n'acheter et de ne lire que des livres dont un pasteur instruit, prudent et connu pour être très-orthodoxe, vous conseillera la lecture. Cependant tout livre qui est bon ne convient pas indifféremment à toute sorte de personnes. Il est sage, souvent même nécessaire, de prendre l'avis de son confesseur avant que d'entreprendre la lecture de quelque livre que ce soit, en matière de religion et de spiritualité.
 
\quarto Mais quand vous voudrez lire, invoquez auparavant le secours d'en-haut par une courte aspiration vers Dieu. Sans la grace intérieure, vous ne tireriez aucun profit surnaturel de quelque lecture que ce fût. Lisez trait à trait pour vous pénétrer de ce qui peut vous être applicable, ou vous instruire. Ne craignez pas de relire plusieurs fois le chapitre qui vous intéresse le plus. Quand vous avez quitté le livre, entretenez-vous intérieurement quelque temps de ce que vous y avez lu. Revenez plusieurs fois dans la journée sur le même objet. Formez toujours quelques résolutions conformes à vos besoins, et que vous mettrez sur-le-champ en pratique. Il est bon encore de commencer vos lectures par la préface, et de suivre méthodiquement jusqu'à la fin, même de relire une seconde et une troisième fois le même livre.
 
\section{Résolutions}
 
\primo Puisque la lecture des bons livres est si utile, même si nécessaire, je ferai ce qui dépendra de moi pour avoir recours tous les jours à ce moyen de salut. 

\secundo Mais je ne lirai point un livre, quelque bon qu'il paraisse, sans avoir préalablement pris l'avis d'un pasteur éclairé et dont la doctrine soit bien connue. 

\tertio Enfin je prendrai les précautions indiquées ci-dessus, afin de tirer de mes lectures un fruit salutaire et solide. 

\prayer{Mon Dieu, présidez à toutes mes lectures par le secours de votre grâce, afin qu'elles me servent à vous connaître, à vous aimer, à me connaître aussi moi-même, et à m'aider dans l'entreprise de mon salut éternel.}
 
\end{document}
